\section[Cote 13 : la question du double réseau et le cas de Tate]{Cote 13 : la question du double réseau et le cas de Tate\protect\verifie[\leanblobmain]{Folder13}}\label{sec:13}

La cote 13, \emph{Groupe de Galois motivique}, est datée de 1965. Aux pages~45 à~49, Grothendieck passe en revue les foncteurs de réalisation des motifs et demande de chacun s'il est pleinement fidèle. La question centrale est celle de la page~47 : sur une extension algébrique $k$ de $\mathbf Q$, le foncteur qui envoie un motif $M$ sur sa réalisation de de Rham $T_{DR}(M)$, sa réalisation de Betti $T_B(M)$ et l'isomorphisme de comparaison $\varphi$ entre elles, sans filtration ni structure de Hodge, est-il pleinement fidèle ? C'est la forme de de Rham--Betti de ce qu'on appelle aujourd'hui la conjecture des périodes de Grothendieck (Bost et Charles 2016, §~2.2 \cite{bostcharles2016}). Elle est ouverte.

\noindent\emph{L'image à garder.} Dans $V \otimes \mathbf C$ cohabitent deux structures rationnelles : le réseau de Betti $V$, rationnel sur $\mathbf Q$, et le réseau de de Rham, rationnel sur $\overline{\mathbf Q}$. La page traduit le second dans le groupe de Galois motivique $G$ : il s'écrit $g(V_{\overline{\mathbf Q}})$ pour un élément $g \in G(\mathbf C)$, la position relative des deux réseaux. Un élément de $V$ invariant par tout $G$ est dans les deux réseaux à la fois ; la question est de savoir si ce sont les seuls. Elle s'écrit
\[
V^G \;\subset\; V \cap g(\overline V) \qquad \text{une égalité pour tout } V \,?
\]
Les rationnels communs aux deux réseaux mesurent les relations algébriques entre périodes : plus les périodes sont transcendantes, moins il y en a.

La page fait deux affirmations sans les démontrer. L'inclusion vaut toujours : la note de la lecture insiste sur son sens. Et le rang un marche : \q{s'il est de dim.\ $1$, c'est une puissance de Tate, et ça marche} (\lot{13}{3}{47}). Les deux sont démontrées ici, dans des modèles choisis par l'édition.

\begin{enonce}[Grothendieck, page 47, l'inclusion $V^G \subset V \cap g(\overline V)$]
Soient un groupe $\Gamma$ agissant sur un ensemble $M$, deux parties $B \subseteq \Lambda$ de $M$ et $g \in \Gamma$. Un élément de $B$ fixé par tout $\Gamma$ appartient à $B \cap g \cdot \Lambda$.
\end{enonce}
Ici $B$ joue le réseau de Betti $V$, $\Lambda$ l'espace $\overline V = V \otimes \overline{\mathbf Q}$ et $g \cdot \Lambda$ le réseau de de Rham. La démonstration tient en une ligne : si $x$ est fixé par tout le groupe, $g^{-1}x = x$, qui est dans $B \subseteq \Lambda$, donc $x \in g \cdot \Lambda$. Elle passe par la description de $V \cap T'(V)$ que donne la page~33 : $x$ appartient à $B \cap g \cdot \Lambda$ si et seulement si $x \in B$ et $g^{-1}x \in \Lambda$ (\texttt{mem\_inter\_smul\_iff}). Lean : \texttt{invariants\_subset\_inter\_smul}.

\begin{enonce}[Grothendieck, page 47, le cas de Tate (forme reconstituée par l'édition)]
Supposons $\pi$ transcendant. Pour tout entier $n$, les rationnels $q$ tels que $q = (2\pi i)^n \alpha$ pour un $\alpha \in \mathbf C$ algébrique sont tous les rationnels si $n = 0$, et le seul $0$ si $n \neq 0$.
\end{enonce}
\emph{Pourquoi c'est le cas de Tate.} Le motif de Tate $\mathbf Q(n)$ a pour réseau de Betti la droite $\mathbf Q \subset \mathbf C$ et pour réseau de de Rham la droite $(2\pi i)^n\, \overline{\mathbf Q}$ : sa seule période est $(2\pi i)^n$. Son groupe de Galois motivique est $\mathbf G_m$, qui agit par la puissance $n$-ième ; les invariants $V^G$ sont donc toute la droite si $n = 0$ et zéro sinon. L'énoncé dit que $V \cap g(\overline V)$, l'ensemble des rationnels qui sont aussi dans la droite de de Rham, est exactement cela : l'égalité $V^G = V \cap g(\overline V)$ vaut en rang un. Cette traduction est le modèle de l'édition ; la page n'écrit ni les deux droites ni le groupe.

\emph{Démonstration.} Pour $n = 0$, tout rationnel $q$ s'écrit $q = 1 \cdot q$, et $q$ est algébrique. Pour $n \neq 0$, supposons $q = (2\pi i)^n \alpha$ avec $q \neq 0$. Alors $\alpha \neq 0$ et $(2\pi i)^n = q\,\alpha^{-1}$ est algébrique, comme produit d'algébriques. Si une puissance entière non nulle de $z$ est algébrique, $z$ l'est : pour un exposant positif, $z$ est racine de $X^n - z^n$ au-dessus d'un corps de nombres ; pour un exposant négatif, on passe à l'inverse (\texttt{isAlgebraic\_of\_zpow}). Donc $2\pi i$ est algébrique, et $\pi = 2\pi i \cdot (2i)^{-1}$ aussi, puisque $i$ est racine de $X^2 + 1$ (\texttt{isAlgebraic\_I}). C'est contraire à l'hypothèse. Lean : \texttt{tate}.

{\small\emph{Ce qui est démontré, et ce qui ne l'est pas.} La transcendance de $\pi$ (Lindemann, 1882) est une hypothèse explicite du théorème : mathlib, à la révision épinglée, a la partie analytique du théorème de Lindemann--Weierstrass et l'irrationalité de $\pi$, mais pas sa transcendance. Le reste est démontré sans hypothèse. Le cas de rang un est donc celui qui \q{marche} parce que la seule période est $(2\pi i)^n$, dont on connaît la transcendance ; dès le rang deux, pour une courbe elliptique, les périodes et quasi-périodes donnent un degré de transcendance attendu égal à~$4$ sans multiplication complexe, ce qui reste ouvert (Chudnovsky établit le cas à multiplication complexe, de degré~$2$). Ne sont pas formalisés : l'identification de $\mathbf Q(n)$ aux deux droites, le groupe $H$ des pages~33 et~34, que la partie \q{Théorie de Hodge} encadre plus bas, et la question générale de la page~47, qui est la conjecture des périodes. On se reportera aux notes de Grothendieck (\lot{13}{2}{33}, \lot{13}{3}{47}).\par}
